\subsection{Bài tập tổng hợp}

\begin{tuluyen}
Sodium hydroxide (NaOH) có nhiều ứng dụng trong sản xuất và đời sống như sản xuất xà phòng, tơ nhân tạo, giấy,... NaOH là một chất rắn hút ẩm mạnh, vì vậy sau khi pha chế nồng độ của dung dịch NaOH cần được xác định lại theo các bước sau:

\textbf{Bước 1:} Cân 1,206 gam mẫu NaOH rắn cho vào cốc thủy tinh, thêm nước cất, khuấy đều cho chất rắn tan hết, sau đó pha thành 500 mL dung dịch (A).

\textbf{Bước 2:} Lấy 10 mL dung dịch (A) cho vào bình tam giác, sau đó nhỏ vào bình 2 giọt chất chỉ thị phenolphthalein, thêm từ từ dung dịch HCl 0,05 M vào hỗn hợp và lắc đều đến khi dung dịch vừa mất màu thì dùng hết 12 mL dung dịch HCl.

Xét các phát biểu sau:
\begin{enumerate}
\item Cho phenolphthalein vào dung dịch (A) thấy dung dịch chuyển sang màu hồng.
\item Nồng độ mol của NaOH trong dung dịch (A) là 0,05 M.
\item Giả sử NaOH chỉ hấp thụ nước trong không khí thì khối lượng nước đã bị hấp thụ vào mẫu NaOH trên là 0,06 gam.
\item NaOH dễ bị nhiệt phân hủy theo phương trình: \ce{2NaOH ->[t^\circ] Na2O + H2O}.
\end{enumerate}
Hãy xác định mỗi phát biểu đúng hay sai và giải thích.
\source{Trích Câu 1 phần đúng--sai, đề tuyển sinh chuyên Hóa Lâm Đồng 2026.}
\dapso{(a) Đúng; (b) Sai; (c) Sai; (d) Sai.}
\begin{loigiaiBT}
Phenolphthalein chuyển hồng trong dung dịch NaOH nên (a) đúng.

Số mol HCl dùng:
\[
n_{\ce{HCl}}=0.012\times0.05=\SI{6.0e-4}{\mole}.
\]
Theo \ce{HCl + NaOH -> NaCl + H2O}, trong \SI{10}{\milli\liter} dung dịch (A) có \SI{6.0e-4}{\mole} NaOH:
\[
C_{\ce{NaOH}}=\frac{6.0\times10^{-4}}{0.010}=\SI{0.060}{\mole\per\liter}.
\]
Vì vậy (b) sai.

Trong \SI{500}{\milli\liter} dung dịch (A):
\[
n_{\ce{NaOH}}=0.060\times0.500=\SI{0.030}{\mole},
\]
\[
m_{\ce{NaOH}}=0.030\times40=\SI{1.20}{\gram}.
\]
Khối lượng nước hấp thụ:
\[
1.206-1.200=\SI{0.006}{\gram}.
\]
Do đó phát biểu (c) với \SI{0.06}{\gram} là sai nếu dữ kiện là 1,206 g như trên.

NaOH bền nhiệt trong điều kiện thông thường của chương trình; phương trình ở (d) không được dùng như một tính chất của NaOH, nên (d) sai.

\textbf{Vậy,} (a) đúng; (b), (c), (d) sai theo dữ kiện đã cho.
\end{loigiaiBT}
\end{tuluyen}

\begin{tuluyen}
Nêu hiện tượng và viết phương trình hóa học của các phản ứng xảy ra khi tiến hành các thí nghiệm sau:
\begin{enumerate}
\item Cho viên kẽm vào dung dịch \ce{KHSO4}.
\item Cho dung dịch \ce{FeCl2} vào dung dịch \ce{AgNO3} dư.
\item Cho dung dịch \ce{Na2CO3} vào dung dịch \ce{FeCl3}.
\end{enumerate}
\source{Trích các ý a), b), d) Câu 1.1, đề tuyển sinh chuyên Hóa Lâm Đồng 2026.}
\dapso{Các phản ứng đều có hiện tượng quan sát được; xem lời giải.}
\begin{loigiaiBT}
Với \ce{KHSO4}, ion \ce{H+} phản ứng với Zn, tạo khí hydrogen:
\[
\ce{Zn + 2KHSO4 -> ZnSO4 + K2SO4 + H2 ^}.
\]
Zn tan dần, có bọt khí không màu.

Khi cho \ce{FeCl2} vào \ce{AgNO3} dư, trước hết tạo kết tủa trắng \ce{AgCl}; đồng thời \ce{Ag+} dư có thể oxi hóa \ce{Fe^{2+}} thành \ce{Fe^{3+}} và tạo Ag kim loại. Đây là trường hợp mở rộng so với phản ứng trao đổi đơn giản.

Với \ce{Na2CO3} và \ce{FeCl3}, phản ứng trong nước không chỉ là trao đổi tạo \ce{Fe2(CO3)3}; xuất hiện kết tủa \ce{Fe(OH)3} và khí \ce{CO2} do sự thủy phân mạnh của ion \ce{Fe^{3+}} và carbonate. Một phương trình tổng quát có thể viết:
\[
\ce{2FeCl3 + 3Na2CO3 + 3H2O -> 2Fe(OH)3 v + 3CO2 ^ + 6NaCl}.
\]

\textbf{Vậy,} cần dựa vào hiện tượng thực tế và bản chất phản ứng trong dung dịch, không chỉ đổi chỗ các ion theo máy móc.
\end{loigiaiBT}
\end{tuluyen}

\begin{tuluyen}
Một nhà máy sử dụng nồi hơi để đun nước giếng khoan có chứa \ce{Mg(HCO3)2} và \ce{Ca(HCO3)2}. Sau một thời gian dài sử dụng thường có cặn rắn Y (chỉ gồm hai muối carbonate) bám dưới đáy nồi hơi.
\begin{enumerate}
\item Xác định hai muối trong Y và viết các phương trình hóa học xảy ra trong nồi hơi để tạo thành hai muối này.
\item Để đánh sạch lớp cặn rắn Y, nhà máy đã sử dụng dung dịch HCl loãng. Viết các phương trình hóa học xảy ra.
\item Trong tự nhiên chúng ta thường bắt gặp hiện tượng những dãy núi đá vôi bị ăn mòn theo thời gian. Hãy dùng kiến thức hóa học để giải thích hiện tượng được đề cập đến trong câu tục ngữ: ``Nước chảy đá mòn''.
\end{enumerate}
\source{Trích nguyên văn Câu 2.1, đề tuyển sinh chuyên Hóa Lâm Đồng 2026.}
\dapso{Y gồm \ce{MgCO3}, \ce{CaCO3}; xem lời giải.}
\begin{loigiaiBT}
Khi đun nóng:
\[
\ce{Ca(HCO3)2 ->[t^\circ] CaCO3 v + CO2 ^ + H2O},
\]
\[
\ce{Mg(HCO3)2 ->[t^\circ] MgCO3 v + CO2 ^ + H2O}.
\]
Vậy Y gồm \ce{CaCO3} và \ce{MgCO3}.

Dùng HCl:
\[
\ce{CaCO3 + 2HCl -> CaCl2 + CO2 ^ + H2O},
\]
\[
\ce{MgCO3 + 2HCl -> MgCl2 + CO2 ^ + H2O}.
\]

Nước tự nhiên có hòa tan \ce{CO2}. \ce{CO2}, nước và \ce{CaCO3} có thể tạo bicarbonate tan:
\[
\ce{CaCO3 + CO2 + H2O -> Ca(HCO3)2}.
\]
Quá trình kéo dài làm đá vôi bị hòa tan dần.

\textbf{Vậy,} cặn trong nồi hơi là carbonate; acid làm tan cặn, còn nước có \ce{CO2} có thể làm đá vôi bị ăn mòn dần theo thời gian.
\end{loigiaiBT}
\end{tuluyen}

\begin{tuluyen}
Một mẫu dung dịch chứa đồng thời \ce{HCl} và \ce{H2SO4}. Để trung hòa hoàn toàn \SI{20.0}{\milli\liter} dung dịch cần \SI{30.0}{\milli\liter} dung dịch \ce{NaOH} \SI{1.00}{\mole\per\liter}. Một mẫu khác cũng có thể tích \SI{20.0}{\milli\liter} được cho tác dụng với \ce{BaCl2} dư, thu được \SI{2.33}{\gram} \ce{BaSO4}. Tính nồng độ mol của mỗi acid. Cho $M_{\ce{BaSO4}}=233$ g/mol.
\source{Bài tổng hợp dạng acid hỗn hợp -- kết tủa -- trung hòa.}
\dapso{$C_{\ce{H2SO4}}=\SI{0.50}{\mole\per\liter}$; $C_{\ce{HCl}}=\SI{0.50}{\mole\per\liter}$.}
\begin{loigiaiBT}
\[
n_{\ce{BaSO4}}=\frac{2.33}{233}=\SI{0.010}{\mole}.
\]
Theo phản ứng tạo \ce{BaSO4}, trong \SI{20.0}{\milli\liter} mẫu có \SI{0.010}{\mole} \ce{H2SO4}:
\[
C_{\ce{H2SO4}}=\frac{0.010}{0.0200}=\SI{0.50}{\mole\per\liter}.
\]
Số mol NaOH dùng:
\[
n_{\ce{NaOH}}=0.0300\times1.00=\SI{0.0300}{\mole}.
\]
\SI{0.010}{\mole} \ce{H2SO4} cần \SI{0.020}{\mole} NaOH; phần còn lại \SI{0.010}{\mole} NaOH trung hòa \ce{HCl}. Do đó:
\[
n_{\ce{HCl}}=\SI{0.010}{\mole},\qquad
C_{\ce{HCl}}=\frac{0.010}{0.0200}=\SI{0.50}{\mole\per\liter}.
\]

\textbf{Vậy,} cả \ce{HCl} và \ce{H2SO4} đều có nồng độ \SI{0.50}{\mole\per\liter}.
\end{loigiaiBT}
\end{tuluyen}

\begin{tuluyen}
Hòa tan hoàn toàn \SI{10.0}{\gram} hỗn hợp \ce{MgO} và \ce{CuO} bằng dung dịch \ce{HCl} vừa đủ. Sau phản ứng thu được \SI{19.625}{\gram} hỗn hợp muối chloride khan. Tính khối lượng mỗi oxide trong hỗn hợp ban đầu. Cho $M_{\ce{MgO}}=40$, $M_{\ce{MgCl2}}=95$, $M_{\ce{CuO}}=80$, $M_{\ce{CuCl2}}=135$ g/mol.
\source{Bài tổng hợp dạng oxide base -- acid -- bảo toàn khối lượng theo PTHH.}
\dapso{\SI{4.0}{\gram} \ce{MgO}; \SI{6.0}{\gram} \ce{CuO}.}
\begin{loigiaiBT}
Gọi số mol \ce{MgO}, \ce{CuO} lần lượt là $x$, $y$.
\[
40x+80y=10.0.
\]
Theo phản ứng với HCl, mỗi mol oxide tạo một mol muối tương ứng:
\[
95x+135y=19.625.
\]
Giải hệ được $x=0.100$ mol, $y=0.075$ mol.
\[
m_{\ce{MgO}}=0.100\times40=\SI{4.0}{\gram},
\]
\[
m_{\ce{CuO}}=0.075\times80=\SI{6.0}{\gram}.
\]

\textbf{Vậy,} hỗn hợp ban đầu chứa \SI{4.0}{\gram} \ce{MgO} và \SI{6.0}{\gram} \ce{CuO}.
\end{loigiaiBT}
\end{tuluyen}
